% Auto-generated by qprov export-latex.
% Define \fact and \provid in your document preamble, e.g.:
%
%   \newcommand{\provid}[1]{\texttt{#1}}
%   \newcommand{\fact}[1]{\textbf{Claim:} #1}
%
% qprov footnotes each fact with its provenance computation id.

\fact{For $[\pi]_q$, the first nonzero Taylor coefficient at $q^k$ for $k > 2$ appears at $k = 10$, with value $c_{10}([\pi]_q) = 1$.}\footnote{Provenance: \provid{bd9390f53345afa6eecf38e9428ed5c9}, recorded 2024-01-15T04:27:55.860412Z.}
% notes: Derived from q_real_truncated('pi', 100) via MGO Prop 1.1.

\fact{The maximum absolute Taylor coefficient of $[\pi]_q$ over the first 50 powers is $\max_{k < 50} |c_k([\pi]_q)| = 51$.}\footnote{Provenance: \provid{bd9390f53345afa6eecf38e9428ed5c9}, recorded 2024-01-15T04:27:55.865434Z.}
% notes: Sanity-checks the MGO bound on coefficient growth.

\fact{For $[\sqrt{2}]_q$, the first negative Taylor coefficient appears at $k = 5$ with value $c_{5}([\sqrt{2}]_q) = -2$.}\footnote{Provenance: \provid{145f920d8f60cd3bee0a08d18a73a2e7}, recorded 2024-01-15T04:27:55.869433Z.}
% notes: Demonstrates that q-real coefficients are not eventually monotone.

\fact{For $[\sqrt{2}]_q$, the maximum absolute coefficient over the first 50 powers is $\max_{k < 50} |c_k([\sqrt{2}]_q)| = 111079093206$.}\footnote{Provenance: \provid{145f920d8f60cd3bee0a08d18a73a2e7}, recorded 2024-01-15T04:27:55.874031Z.}

\fact{For the golden ratio $\varphi$, the maximum absolute coefficient of $[\varphi]_q$ over the first 25 powers is $\max_{k < 25} |c_k([\varphi]_q)| = 37150472$, which exhibits the expected exponential growth at rate $\varphi$.}\footnote{Provenance: \provid{ed093c18aee16baeb02cbfd9bb62227c}, recorded 2024-01-15T04:27:55.878061Z.}

\fact{For Euler's $e$, the number of zero Taylor coefficients of $[e]_q$ among the first 50 powers is $2$.}\footnote{Provenance: \provid{5662672805bfb479ade1ea6f3d5dbfe0}, recorded 2024-01-15T04:27:55.883678Z.}

\fact{For the rational $[3/2]_q = \frac{1+q+q^2}{1+q}$, the first negative Taylor coefficient appears at $k = 3$, consistent with the geometric expansion of $1/(1+q)$.}\footnote{Provenance: \provid{7eb093735dbfa9a26fd521426bb633c2}, recorded 2024-01-15T04:27:55.888256Z.}
